%% ch13 --- Zestaw inzyniera AI: pydantic, httpx i ksztalt SDK
%%
%% Napisane we wrzesniu 2026. Kazde zdanie o SDK w tym rozdziale wypisuje
%% code/ch13/sdk_shape.py z zainstalowanego pakietu, a kazda liczba pochodzi
%% z code/measure/e08_validation.py. Nic tutaj nie zostalo napisane z
%% pamieci; tam, gdzie opis rozdzialu okazal sie bledny, mowi o tym
%% CLAUDE.md w sekcji *Resolved questions*.

\chapter{Zestaw inżyniera AI: pydantic, httpx i kształt SDK}
\label{ch:aitoolkit}

\section{Trzy rzeczy w jednym płaszczu}
\label{sec:ai-three-things}
\index{SDK!z czego się składa}

Dodałeś do projektu \pkg{openai} albo \pkg{anthropic} i edytor podpowiada ci
sto nazw. Przychodząc z .NET, rozsądnie jest spodziewać się SDK w stylu
dostawcy: wygenerowanego klienta z własnym modelem obiektowym, własnym typem
opcji, własnym łańcuchem handlerów i warstwą abstrakcji między tobą a
drutem, z którą prędzej czy później przyjdzie ci walczyć.

To oczekiwanie jest błędne, i to w stronę, która ułatwia ci życie. Zanim
pójdziesz dalej, odpowiedz: wywołanie SDK oddaje ci obiekt. Nie \dash{} co
jest w środku, tylko jakiej \emph{klasy} jest to instancja?

Odpowiedź brzmi: modelu pydantic. To, co go pobrało, jest klientem httpx.
Strumieniową postacią tego samego wywołania jest iterator. I to wszystko:
trzy kawałki Pythona, które już umiesz czytać, złożone tak, jak sam byś je
złożył.

\mermaidfig{sdk-anatomy}{Z czego zbudowany jest każdy SDK modelu w Pythonie.
Każda z tych trzech rzeczy to biblioteka, którą da się przeczytać, a SDK
jest okablowaniem między nimi.}{trzy części SDK modelu}

\begin{csbox}
\code{HttpClient} plus \code{System.Text.Json} plus
\code{IAsyncEnumerable<T>}, a SDK jest kodem pomiędzy nimi. Czego Python nie
ma, to \code{IHttpClientFactory}: żaden kontener nie podaje ci klienta z
puli, więc decyzja o czasie jego życia należy do ciebie, a zapada w
\S\ref{sec:ai-which-httpx}.
\end{csbox}

Ten rozdział kończy się na jednym wywołaniu modelu i jednym ustrukturyzowanym
wyjściu. Agenci, narzędzia, middleware i grafy to inna książka \dash{} są
tematem tomu towarzyszącego, a powtarzanie ich tutaj dałoby ci dwa opisy
jednej rzeczy, które się rozjadą. Dostajesz za to podłogę, na której one
stoją, czyli tę część, która nie zmienia się razem z frameworkiem.

\section{Ustrukturyzowane wyjście to klasa, którą już napisałeś}
\label{sec:ai-structured}
\index{pydantic!ustrukturyzowane wyjście}

Model zwraca tekst. \emph{Ustrukturyzowane wyjście} to ten tekst sparsowany
do typu, który zadeklarowałeś, a w Pythonie tym typem jest model pydantic.
Pydantic spotkałeś na granicy usługi w rozdziale~\ref{ch:services}; to ten
sam mechanizm obrócony w drugą stronę \dash{} ku odpowiedzi, nie ku żądaniu.

\pyregion{code/ch13/structured.py}{models}{Dwie deklaracje jednego kształtu.
Druga różni się o jedną linię, a \S\ref{sec:ai-what-costs} mierzy, ile ta
linia kosztuje.}{lst:ch13-models}

Całą robotę wykonują dwa wywołania i razem są one tym, czym jest pomocnik od
ustrukturyzowanych wyjść w SDK.

\pyregion{code/ch13/structured.py}{schema}{W stronę drutu: model staje się
schematem JSON, a SDK dokłada do niego to, czego chce
dostawca.}{lst:ch13-schema}

\mermaidfig{schema-round-trip}{Jedna klasa użyta dwa razy: raz, żeby
powiedzieć, jak ma wyglądać odpowiedź, i raz, żeby sprawdzić, że tak
wyglądała.}{ustrukturyzowane wyjście, tam i z powrotem}

W stronę drutu \api{model\_json\_schema()} zamienia twoją klasę na schemat
JSON, który niesie żądanie. W drugą stronę
\api{model\_validate\_json()}\index{model\_validate\_json@\texttt{model\_validate\_json}}
zamienia tekst odpowiedzi w instancję tej klasy. Wkładem SDK do schematu są
dwa klucze \dash{} nazwa oraz \code{additionalProperties: false} dla
dostawcy, który obiecuje zwrócić twoje pola i nic poza nimi. Cała reszta na
drucie należy do pydantica.

Warto się przy tym zatrzymać, bo to odpowiedź na pytanie, które inaczej
wygląda głęboko. Nie ma rejestru, nie ma kroku generowania kodu i nie ma
warstwy mapującej. Kształt jest zapisany w jednym miejscu, w klasie, i obie
połowy tej podróży czytają go stamtąd.

\section{Co się przeciska}
\label{sec:ai-what-gets-through}
\index{pydantic!tryb ścisły i luźny}

Teraz część, w której twoje nawyki z .NET aktywnie przeszkadzają.
Zadeklarowałeś \code{severity: int}. Dostawca, albo proxy przed nim, przysyła
\code{"2"} \dash{} właściwą wartość, jako napis. Zanim pójdziesz dalej: co
zrobi z tym \api{model\_validate\_json()}?

\pyregion{code/ch13/structured.py}{validate}{Ten sam ładunek, dwa razy. Jedno
z tych zachowań jest tym, co zrobiłby
\code{JsonSerializer.Deserialize}.}{lst:ch13-validate}

\transcript{ch13-structured}

Konwertuje. \code{"2"} staje się \num{2}, \code{"true"} staje się
\code{True}, i nikt nigdzie nie zostaje o tym poinformowany. To domyślne
zachowanie pydantica i nazywa się trybem \emph{luźnym}.

\begin{trapbox}
\code{System.Text.Json} odmawia przyjęcia napisu JSON dla właściwości
liczbowej, dopóki nie zgodzisz się na to przez \code{NumberHandling}.
Pydantic działa odwrotnie: konwersja jest domyślna, a odmowa jest opcją.
Odruch mówiący, że \emph{deserializacja waliduje}, obowiązuje więc w C\# i
nie obowiązuje tutaj \dash{} granica napisana z tego odruchu po cichu
przyjmuje zmianę kontraktu, a pierwszą rzeczą, jakiej się o niej dowiesz,
jest zła liczba gdzieś dalej.

\end{trapbox}

\code{strict=True} w konfiguracji modelu odkupuje zachowanie, którego się
spodziewałeś. Użyteczne pytanie nie brzmi, który tryb jest słuszny w
oderwaniu od wszystkiego, tylko co każdy z nich przepuszcza \dash{} a to da
się policzyć.

\pyregion{code/ch13/what_survives.py}{plain}{Trzecia strategia: typ, który
inżynier C\# pisze najpierw.}{lst:ch13-plain}

\pyfile{code/ch13/what_survives.py}{Siedem odpowiedzi, które są poprawnym
JSON-em i są złe. Eksperyment E8 importuje ten zestaw, zamiast go kopiować,
więc ta tabela i ten pomiar nie mogą się rozjechać.}{lst:ch13-survives}

\transcript{ch13-survives}

Przeczytaj najpierw ostatni wiersz: z \val{e08.corpus} zepsutych odpowiedzi
tryb luźny odrzuca \val{e08.reject.lax}, ścisły odrzuca
\val{e08.reject.strict}, a dataclass odrzuca \val{e08.reject.plain}.
Dataclass przepuszcza \val{e08.silent.plain} z nich, każdą złą, ponieważ
dataclass generuje \api{\_\_init\_\_} i nie sprawdza w czasie działania ani
jednej ze swoich adnotacji.

Przeczytaj teraz ostatnią linię tabeli, bo to ona jest warta ceny tego
rozdziału. Jedyny ładunek, który dataclass odrzuca, to ten, z którym było
\emph{wszystko w porządku}: dostawca dodał pole. Oba SDK spodziewają się, że
tak będzie \dash{} ich własne modele są skonfigurowane tak, by zachowywać
nieznane pola, a nie porzucać je, co wypisuje \S\ref{sec:ai-which-httpx}
\dash{} więc naiwna strategia notuje jedną odmowę i jest to fałszywy alarm.

\section{Ile to kosztuje, zmierzone}
\label{sec:ai-what-costs}
\index{pomiar!E8, koszt walidacji}

Oczywiste następne pytanie brzmi, ile to bezpieczeństwo kosztuje, a oczywiste
oczekiwanie to wymiana: walidacja jest pracą, praca zajmuje czas, wybierz
punkt na krzywej. Eksperyment E8 mierzy czasy trzech strategii na zestawie
powyżej i to oczekiwanie tego nie przeżywa.

Dwa wyniki, na maszynie, na której to powstawało. \textbf{Tryb ścisły nie
kosztuje nic ponad luźny} \dash{} oba wyszły w granicach kilku procent od
siebie, raz w jedną, raz w drugą stronę. Oraz: \textbf{sparsowanie i
walidacja przez pydantica bije sam parser biblioteki standardowej}: na
odpowiedzi o \val{e08.fields.wide} polach samo \api{json.loads} zajęło
ponad \val{e08.json.factor} razy więcej czasu, niż
\api{model\_validate\_json} potrzebowało, żeby sparsować te same bajty
\emph{i} sprawdzić każde pole względem typu.

Mechanizm przestaje być tajemniczy, gdy się go wypowie. \api{json.loads}
buduje pythonowy słownik z każdego klucza i każdej wartości, którą przysłał
dostawca, a potem większość tego wyrzucasz. Rdzeń pydantica parsuje bajty w
skompilowanym kodzie prosto do pól modelu i nigdy nie tworzy pośredniego
słownika. Im szersza odpowiedź, tym bardziej odjeżdża \dash{} co jest
odwrotnością kształtu, który narysowałbyś, myśląc o walidacji jak o podatku
nałożonym na parsowanie.

\begin{note}
Żadnej z tych dwóch wartości nie ma na tej stronie i jest to celowe.
Mikrosekunda jest własnością maszyny, więc nie da się jej zapisać w książce,
która sama przelicza swoje pomiary i przewraca się, gdy któryś się ruszy. E8
zapisuje w zamian parę sufitów, dobranych z zapasem i
\emph{sprawdzanych przy każdym uruchomieniu}: tryb ścisły może kosztować
najwyżej \val{e08.strict.overhead.pct} procent więcej niż luźny, a samo
\api{json.loads} musi być co najmniej \val{e08.json.factor} razy wolniejsze.
Każdy z nich brany jest jako najlepszy z \val{e08.repeats} przebiegów po
\val{e08.iterations} wywołań, bo najszybszy przebieg jest tym najmniej
zaburzonym przez planistę, podczas gdy średnia jest średnią z zaburzeń.
Jeśli późniejszy pydantic złamie którykolwiek z nich, build robi się
czerwony i ta sekcja jest błędna tego samego dnia. Liczby w
\S\ref{sec:ai-what-gets-through} nie są sufitami: są własnością kodu i są
identyczne na każdej maszynie.
\end{note}

Wymiana, od której zaczęła się ta sekcja, nie istnieje, a rada jest wyjątkowo
prosta. Włącz tryb ścisły na każdej granicy, którą przekracza odpowiedź
modelu. To jedna linia, odrzuca dwie dodatkowe klasy zepsutych odpowiedzi i
jest za darmo.

\section{Który httpx}
\label{sec:ai-which-httpx}
\index{httpx!która dystrybucja}

Wszystko powyżej dotyczy nogi pydanticowej. W nodze transportowej jest
pułapka, przed którą nie ostrzeże cię żadna dokumentacja, a jedynym sposobem,
żeby ją znaleźć, jest przeczytanie zainstalowanego pakietu.

Przeczytaj go więc. Klient w środku SDK jest klientem httpx, a ta książka
przypina \pkg{httpx}. Zanim pójdziesz dalej: co zwróci
\api{isinstance(client.\_client, httpx.Client)}?

\pyfile{code/ch13/sdk_shape.py}{Każde zdanie, jakie ten rozdział mówi o SDK,
wypisane z zainstalowanego pakietu. Asercje na dole oznaczają, że CI padnie w
dniu, w którym któreś z nich przestanie być prawdziwe.}{lst:ch13-sdk-shape}

\transcript{ch13-sdk-shape}

\pkg{httpx} i \pkg{httpx2} to dwie dystrybucje. Ta książka przypina
pierwszą, w wersji~\httpxver, bo to z nią pisze się klienta; oba SDK zależą
od drugiej, a klient w ich środku jest klientem \pkg{httpx2}. Mają wspólne
pochodzenie i wspólną pisownię, i nie są tym samym pakietem, więc
\api{isinstance(client.\_client, httpx.Client)} daje \code{False} na kliencie
SDK, który jest przecież oczywistym klientem httpx.

\begin{warning}
Oba SDK nie zgadzają się co do tego, czy podanie im niewłaściwego klienta
jest błędem. Zmierzone, nie wywnioskowane: \pkg{anthropic} odrzuca
\api{httpx.Client} błędem \code{TypeError}, który nazywa oba pakiety, a
\pkg{openai} przyjmuje go i działa \dash{} przynajmniej na zwykłym żądaniu, i
tylko tyle zostało sprawdzone. \enquote{U mnie zadziałało} nie jest zatem
dowodem, że klient, którego podałeś, jest tym, którego SDK oczekuje.
\end{warning}

Z tego listingu wynikają jeszcze dwie rzeczy. Własne modele SDK są
skonfigurowane tak, by \emph{zachowywać} pola, których nikt nie
zadeklarował, zamiast je ignorować \dash{} i to właśnie sprawia, że dodanie
pola przez dostawcę jest nie-wydarzeniem. Domyślne limity czasu natomiast nie
są czymś, co chcesz odziedziczyć: limit odczytu, który wypisuje listing, to
dziesięć minut, co jest rozsądnym ustawieniem dla dostawcy strumieniującego długą odpowiedź i
fatalnym dla żądania siedzącego w handlerze HTTP, na który ktoś czeka.
Rozdział~\ref{ch:services} daje temu żądaniu budżet; SDK nie da mu go za
ciebie.

Drugą decyzją jest czas życia klienta. Klient httpx trzyma pulę połączeń,
więc jeden na proces jest właściwy, a jeden na żądanie wyrzuca tę pulę za
każdym razem \dash{} ten sam argument i ten sam tryb awarii, co przy czasie
życia \api{AsyncClient} w rozdziale~\ref{ch:asyncio}.

\section{Tokeny są wielkościami, a pieniądze nie są floatem}
\label{sec:ai-tokens}
\index{koszt!tokeny jako wielkość}

Każda odpowiedź niesie liczby, które decydują o rachunku, w polu
\api{usage}, będącym kolejnym modelem pydantic. Usługa, która ich nie
sumuje, dowiaduje się, ile wydała, z faktury \dash{} a wtedy przebieg, który
te pieniądze wydał, ma już dwa tygodnie.

\begin{csbox}
Cena to pieniądze, więc jest typu \api{Decimal}, a nigdy \code{float},
dokładnie z tego powodu, dla którego w C\# jest to \code{decimal}, a nigdy
\code{double}. Różnica polega na tym, że kompilator C\# każe ci powiedzieć,
co miałeś na myśli, a Python nie: nieotypowany literał jest floatem,
\code{0.1 + 0.2} nie jest równe \code{0.3} także tutaj, a miesiąc kosztów
pojedynczych wywołań zsumowany na floatach dryfuje. \api{Decimal} bierze
też swoją precyzję z kontekstu, a nie z typu, więc kwantuj świadomie,
zamiast ufać wartości domyślnej.
\end{csbox}

Ćwiczenie~\ref{ex:ch13-cost} to ta arytmetyka, zrobiona raz i porządnie.

\section{Notatniki kontra skrypty}
\label{sec:ai-notebooks}
\index{notatniki!i dlaczego nie ma ich w tej książce}

Każdy tutorial, jaki spotkasz do tego materiału, jest notatnikiem, a nic w
tej książce nim nie jest. To stanowisko, więc oto argument za nim.

Notatnik nie jest plikiem, który się wykonuje. Jest REPL-em, który pamięta
każdą komórkę, jaką uruchomiłeś, w kolejności, w jakiej je uruchamiałeś
\dash{} a nie w tej, w jakiej są na ekranie. Usuń komórkę, a nazwa, którą
zdefiniowała, nadal jest związana. Uruchom je od góry do dołu na świeżym
jądrze, a dostaniesz inny wynik niż ten na ekranie, i to w tej różnicy
przepada popołudnie. Nic w tej książce nie mogłoby być notatnikiem, bo każdy
listing tutaj wykonuje się na przypiętych wersjach przy każdym pushu, a
\enquote{u mnie w jądrze działało} nie jest czymś, co CI potrafi sprawdzić.

Notatnik jest właściwy do pracy, do której powstał: oglądania danych,
próbowania czegoś, trzymania wykresu obok kodu, który go narysował. Używaj
go. Zasada warta zapamiętania jest węższa niż \emph{unikaj notatników}
\dash{} brzmi ona, że nic, co importuje usługa, nie mieszka w notatniku. W
momencie, w którym funkcję ma wywołać coś innego niż ty, jej miejsce jest w
module, z testem obok.

\section{Jedno wywołanie, otypowane z obu stron}
\label{sec:ai-call-model}

Oto wszystko powyżej, złożone w najmniejszą użyteczną rzecz: jedno wywołanie
z typem wchodzącym, typem wychodzącym, budżetem czasu i zapisem tego, co się
wydarzyło.

\pyfile{code/ch13/call_model.py}{Nagroda. Nic tutaj nie jest specyficzne dla
dostawcy i wchodzi do usługi z rozdziału~\ref{ch:services} tak, jak
stoi.}{lst:ch13-call-model}

\transcript{ch13-call-model}

Trzy szczegóły zasługują na swoje miejsce. Limit czasu jest argumentem z
wartością domyślną liczoną w sekundach, a nie dziesięcioma minutami z SDK.
Awaria jest zgłaszana jako jeden własny typ wyjątku, więc wywołujący obsługuje
\emph{model nie odpowiedział}, zamiast uczyć się taksonomii błędów dostawcy
\dash{} i nazywa się \code{...Error}, a nie \code{...Exception}, co jest
konwencją Pythona i czego pilnuje linter. A odpowiedź jest sprawdzana pod
kątem tego, czy w ogóle się sparsowała, zanim zostanie zwrócona, bo
\api{output\_parsed} jest opcjonalne, a wywołujący, który to zignoruje,
dostaje \code{None} dwie ramki stosu od przyczyny.

\mermaidfig{two-boundaries}{Dwie krawędzie usługi i ten sam mechanizm na
obu.}{walidacja na obu granicach}

Ten listing wykonuje się w CI przy każdym pushu, bez dostawcy i bez klucza
API. Lokalny serwer odpowiadający w kształcie, który SDK parsuje, jest dla
SDK nie do odróżnienia od dostawcy \dash{} i to właśnie sprawia, że rozdział
o SDK modeli daje się w ogóle testować.

\begin{projectbox}
\textbf{Etap 02 \dash{} zapis wywołania modelu.} Model śladu z
rozdziału~\ref{ch:testing} zostawił jedną rzecz otwartą: co wywołanie modelu
wkłada do pola \code{tags} zdarzenia. Ten rozdział to rozstrzyga. Wywołanie
to \emph{dwa} zdarzenia, a nie jedno, bo wywołanie, które nigdy nie wróciło,
jest przypadkiem, który najbardziej chcesz umieć znaleźć, a pojedyncze
zdarzenie zapisane po odpowiedzi nie potrafi go przedstawić. Pole
\code{tags} niesie liczby tokenów i to, czy odpowiedź się sparsowała; nie
niesie promptu ani odpowiedzi modelu, bo ślad jest przechowywany, przesyłany
i czytany przez ludzi, których nie było w pokoju.
\end{projectbox}

\section{Ćwiczenia}
\label{sec:ai-exercises}

\begin{exercise}{e13_01_strict_boundary}{Granica, która nie zgaduje}
\code{Reply} konwertuje, więc napis \code{"2"} przychodzi jako liczba
całkowita i nikt nie zostaje o tym poinformowany. Spraw, żeby parsowanie
odrzucało wszystko, czego typy nie są dokładnie tym, co deklaruje klasa, i
zwracało zamiast tego \code{None}. Nie pisz tych sprawdzeń ręcznie.
\end{exercise}

\begin{exercise}{e13_02_wire_schema}{Schemat, który dostawca przyjmie}
Dostawca, który obiecuje zwrócić twoje pola i nic poza nimi, chce
\code{additionalProperties: false} na każdym obiekcie w schemacie. Model
zagnieżdżony jest drugim obiektem, a pydantic nie zagnieżdża go tam, gdzie
byś się spodziewał \dash{} wyciąga go na zewnątrz. Domknij je wszystkie.
\end{exercise}

\begin{exercise}{e13_03_what_it_cost}{Ile kosztowało jedno wywołanie}
\label{ex:ch13-cost}
Wyceń odpowiedź na podstawie jej dwóch liczb tokenów i dwóch stawek
dostawcy. Wynikiem są pieniądze, więc nie jest floatem; i jest
skwantowany, żeby miesiąc takich wyników sumował się dwa razy do tej samej
liczby.
\end{exercise}

\begin{exercise}{e13_04_read_the_trace}{Dlaczego wywołanie to dwa zdarzenia}
Napisz dwa pytania, dla których istnieje ślad: które modele zostały wywołane
i nigdy nie odpowiedziały oraz ile wydał cały przebieg. Żadne z nich nie może
zakładać, że ślad jest poprawnie zbudowany \dash{} po to się go ma.
\end{exercise}

\section{Gdzie to się kończy}
\label{sec:ai-stops}

Potrafisz już czytać SDK, a nie tylko je wywoływać. Kiedy pojawi się
następne, z inną nazwą zasobu i inną kopertą, rozstrzygają je trzy pytania:
jaki model pydantic wraca, który klient HTTP jest pod spodem i co oddaje
strumieniowa postać wywołania. Odpowiedzi są w zainstalowanym pakiecie i
nigdzie indziej \dash{} nie we wpisie na blogu, nie we wspomnieniach modelu i
nie zawsze w dokumentacji.

Celowo nie ma tu wszystkiego, co jest zbudowane \emph{na} tym: wywoływania
narzędzi, pętli agenta, ponowień z osądem w środku, grafów i stanu, który im
towarzyszy. To temat tomu towarzyszącego i zaczyna się tam, gdzie kończy się
ten rozdział. Rozdział~\ref{ch:traceassert} idzie drugą drogą: kończy
przyrząd, który ten rozdział nauczył zapisywać wywołanie, i wysyła go w
świat.
